% AWS Cloud Architecture and Serverless Implementation

\section{Lambda Function and API Gateway Integration}

\subsection{Serverless Static Website Architecture}

The internship cloud projects demonstrate a modern serverless stack:

\begin{itemize}
    \item \textbf{Frontend:} React Static Website (built and deployed to S3)
    \item \textbf{API:} AWS API Gateway (REST endpoints)
    \item \textbf{Backend:} Lambda Function (Node.js runtime)
    \item \textbf{Database:} DynamoDB (NoSQL document store)
    \item \textbf{Deployment:} CI/CD pipeline using GitHub Actions and Elastic Beanstalk
\end{itemize}

\subsection{Lambda Function for DynamoDB Record Retrieval}

\begin{lstlisting}[language=JavaScript, caption=Lambda Handler for GET Records]
const AWS = require('aws-sdk');
const dynamodb = new AWS.DynamoDB.DocumentClient();

/**
 * Lambda handler to retrieve records from DynamoDB
 * Triggered by API Gateway GET request
 */
exports.handler = async (event) => {
  console.log('Received event:', event);

  try {
    // Extract table name and optional filter from event
    const tableName = process.env.DYNAMODB_TABLE || 'Records';
    const queryParams = event.queryStringParameters || {};

    let params = {
      TableName: tableName
    };

    // Add optional KeyConditionExpression if querying by partition key
    if (queryParams.id) {
      params.KeyConditionExpression = 'id = :id';
      params.ExpressionAttributeValues = {
        ':id': queryParams.id
      };
    } else {
      // Default: Scan all items with optional filter
      if (queryParams.filter) {
        params.FilterExpression = `#status = :status`;
        params.ExpressionAttributeNames = {
          '#status': 'status'
        };
        params.ExpressionAttributeValues = {
          ':status': queryParams.filter
        };
      }
    }

    // Add pagination
    if (queryParams.limit) {
      params.Limit = parseInt(queryParams.limit, 10);
    }

    // Execute DynamoDB query/scan
    const result = await dynamodb.query(params).promise()
      .catch(() => dynamodb.scan(params).promise());

    return {
      statusCode: 200,
      headers: {
        'Content-Type': 'application/json',
        'Access-Control-Allow-Origin': '*'
      },
      body: JSON.stringify({
        success: true,
        count: result.Items.length,
        items: result.Items,
        lastEvaluatedKey: result.LastEvaluatedKey
      })
    };

  } catch (error) {
    console.error('Error retrieving records:', error);

    return {
      statusCode: 500,
      headers: {
        'Content-Type': 'application/json',
        'Access-Control-Allow-Origin': '*'
      },
      body: JSON.stringify({
        success: false,
        message: 'Failed to retrieve records',
        error: error.message
      })
    };
  }
};
\end{lstlisting}

\subsection{API Gateway Integration Configuration}

\begin{lstlisting}[language=json, caption=API Gateway REST API Setup]
{
  "swagger": "2.0",
  "info": {
    "title": "Internship Records API",
    "version": "1.0.0"
  },
  "paths": {
    "/records": {
      "get": {
        "summary": "Retrieve records from DynamoDB",
        "consumes": ["application/json"],
        "produces": ["application/json"],
        "parameters": [
          {
            "name": "id",
            "in": "query",
            "description": "Record ID",
            "type": "string"
          },
          {
            "name": "filter",
            "in": "query",
            "description": "Filter by status",
            "type": "string"
          },
          {
            "name": "limit",
            "in": "query",
            "description": "Maximum records to return",
            "type": "integer"
          }
        ],
        "x-amazon-apigateway-integration": {
          "type": "aws_proxy",
          "httpMethod": "POST",
          "uri": "arn:aws:apigateway:us-east-1:lambda:path/2015-03-31/functions/arn:aws:lambda:us-east-1:ACCOUNT_ID:function:GetRecords/invocations",
          "responses": {
            "default": {
              "statusCode": "200"
            }
          }
        }
      }
    }
  }
}
\end{lstlisting}

\section{CI/CD Pipeline with Elastic Beanstalk}

\subsection{GitHub Actions Automated Workflow}

\begin{lstlisting}[language=yaml, caption=GitHub Actions CI/CD Pipeline]
name: Deploy to Elastic Beanstalk

on:
  push:
    branches:
      - main
      - develop
  pull_request:
    branches:
      - main

jobs:
  test-and-build:
    runs-on: ubuntu-latest
    
    strategy:
      matrix:
        node-version: [18.x, 20.x]

    steps:
      # Checkout code
      - uses: actions/checkout@v3

      # Setup Node.js
      - name: Use Node.js ${{ matrix.node-version }}
        uses: actions/setup-node@v3
        with:
          node-version: ${{ matrix.node-version }}
          cache: 'npm'

      # Install dependencies
      - name: Install dependencies
        run: npm ci

      # Run linting
      - name: Run ESLint
        run: npm run lint --if-present

      # Run tests
      - name: Run tests
        run: npm test --if-present

      # Build application
      - name: Build application
        run: npm run build --if-present

      # Build React frontend (if applicable)
      - name: Build React app
        run: |
          cd frontend
          npm ci
          npm run build
        continue-on-error: true

  deploy:
    needs: test-and-build
    runs-on: ubuntu-latest
    if: github.ref == 'refs/heads/main' && github.event_name == 'push'

    steps:
      - uses: actions/checkout@v3

      # Configure AWS credentials
      - name: Configure AWS credentials
        uses: aws-actions/configure-aws-credentials@v2
        with:
          aws-access-key-id: ${{ secrets.AWS_ACCESS_KEY_ID }}
          aws-secret-access-key: ${{ secrets.AWS_SECRET_ACCESS_KEY }}
          aws-region: us-east-1

      # Install EB CLI
      - name: Install Elastic Beanstalk CLI
        run: |
          python -m pip install --upgrade pip
          pip install awsebcli

      # Deploy to Elastic Beanstalk
      - name: Deploy to Elastic Beanstalk
        run: |
          eb init -p "Node.js 18 running on 64bit Amazon Linux 2" \
            --region us-east-1 \
            --application-name internship-app
          eb create internship-env --instance-type t3.micro || \
          eb deploy internship-env

      # Verify deployment
      - name: Verify deployment
        run: |
          HEALTH=$(eb health --refresh | grep "Overall Health")
          echo "$HEALTH"

      # Send deployment notification
      - name: Notify deployment
        if: always()
        uses: 8398a7/action-slack@v3
        with:
          status: ${{ job.status }}
          text: 'Deployment ${{ job.status }} at ${{ github.server_url }}/${{ github.repository }}/actions/runs/${{ github.run_id }}'
          webhook_url: ${{ secrets.SLACK_WEBHOOK }}
        continue-on-error: true
\end{lstlisting}

\subsection{Elastic Beanstalk Configuration}

\begin{lstlisting}[language=yaml, caption=.ebextensions/app.config]
option_settings:
  # Application settings
  aws:elasticbeanstalk:application:environment:
    NODE_ENV: production
    LOG_LEVEL: info
    DYNAMODB_TABLE: InternshipRecords
  
  # Auto-scaling configuration
  aws:autoscaling:asg:
    MinSize: 1
    MaxSize: 4
  
  aws:autoscaling:trigger:
    MeasureName: CPUUtilization
    Statistic: Average
    Unit: Percent
    UpperThreshold: 70
    LowerThreshold: 30
  
  # Load balancer settings
  aws:elasticbeanstalk:environment:process:default:
    HealthCheckPath: /health
    HealthyThresholdCount: 3
    UnhealthyThresholdCount: 5
  
  # Environment
  aws:elasticbeanstalk:cloudwatch:logs:
    StreamLogs: true
    DeleteOnTerminate: false

files:
  "/opt/elasticbeanstalk/tasks/bundlelogs.d/01_application.conf":
    content: |
      /var/app/current/logs/*.log
      /var/app/current/*.log

commands:
  01_npm_install:
    command: "npm ci --production"
    cwd: /var/app/staging/
  02_npm_build:
    command: "npm run build --if-present"
    cwd: /var/app/staging/
\end{lstlisting}

\newpage
